\documentclass[12pt,a4paper]{article}
\usepackage[utf8]{inputenc}
\usepackage[french]{babel}
\usepackage[T1]{fontenc}
\usepackage{amsmath, amsfonts, amssymb}
\usepackage{graphicx}
\usepackage{geometry}
\usepackage{xcolor}
\usepackage{titlesec}
\usepackage{hyperref}
\usepackage{listings}
\usepackage{booktabs}
\usepackage{fancyhdr}

\geometry{margin=2.5cm}

\definecolor{codeblue}{RGB}{0,70,160}
\definecolor{codegray}{RGB}{245,245,245}

\lstset{
    basicstyle=\ttfamily\small,
    breaklines=true,
    frame=single,
    backgroundcolor=\color{codegray},
    keywordstyle=\bfseries\color{codeblue},
    commentstyle=\itshape\color{gray},
    showstringspaces=false,
    numbers=left,
    numberstyle=\tiny\color{gray},
    numbersep=8pt
}

\pagestyle{fancy}
\fancyhf{}
\rhead{NEUSSI Patrice -- 24P820}
\lhead{Realite Augmentee -- TP4}
\rfoot{\thepage}

\titleformat{\section}{\Large\bfseries\color{codeblue}}{}{0em}{\thesection\quad}
\titleformat{\subsection}{\large\bfseries}{}{0em}{\thesubsection\quad}
\hypersetup{colorlinks=true, linkcolor=codeblue, urlcolor=codeblue}

\begin{document}

\begin{titlepage}
    \centering
    \vspace*{2cm}
    {\huge\bfseries Rapport de Travaux Pratiques\par}
    \vspace{1cm}
    {\LARGE\bfseries\color{codeblue} TP4 -- Quaternions et Rotations\par}
    \vspace{1.5cm}
    \rule{\linewidth}{0.5pt}\vspace{0.5cm}
    {\large Cours : Realite Augmentee FROM SCRATCH\par}
    {\large Module 04 -- NkMath\par}
    \vspace{0.5cm}\rule{\linewidth}{0.5pt}
    \vspace{2cm}
    {\large\textbf{Etudiant :} NEUSSI NJIETCHEU PATRICE EUGENE\par}
    {\large\textbf{Matricule :} 24P820\par}
    \vspace{2cm}
    {\large Repository GitHub :\par}
    \href{https://github.com/neussi/RA_TP4_Quaternions_et_Rotations}{https://github.com/neussi/RA\_TP4\_Quaternions\_et\_Rotations}
    \vfill
    {\large \today\par}
\end{titlepage}

\tableofcontents
\newpage

\section{Introduction}

Ce rapport presente le module \textbf{Quaternions et Rotations} du cours \textit{Realite
Augmentee FROM SCRATCH}. Les quaternions evitent les problemes de gimbal lock des angles
d'Euler et permettent une interpolation spherique fluide (SLERP) indispensable en AR.

\section{Concepts Theoriques}

\subsection{Definition d'un Quaternion}

Un quaternion est un element de $\mathbb{H}$ :
\[
    q = w + xi + yj + zk \qquad \text{avec} \quad i^2=j^2=k^2=-1,\; ij=k
\]

Le quaternion unitaire $||q||=1$ decrit une rotation pure d'angle $\theta$ autour de l'axe $\hat{n}$ :
\[
    q = \cos\frac{\theta}{2} + \hat{n}\sin\frac{\theta}{2}
\]

\subsection{Operations Fondamentales}

\begin{align}
    q^* &= w - xi - yj - zk \quad \text{(conjugue)} \\
    q^{-1} &= \frac{q^*}{||q||^2} \quad \text{(inverse, } q^{-1}=q^* \text{ si unitaire)} \\
    q_1 q_2 &: \text{produit de Hamilton} \\
    q p q^* &: \text{rotation du vecteur } p
\end{align}

\subsection{SLERP -- Interpolation Spherique}

\[
    \text{Slerp}(q_1, q_2, t) = q_1 \left(\frac{\sin((1-t)\Omega)}{\sin\Omega}\right)
    + q_2 \left(\frac{\sin(t\Omega)}{\sin\Omega}\right)
\]
avec $\Omega = \arccos(q_1 \cdot q_2)$.

\section{Realisation Technique}

\begin{lstlisting}[language=C++, caption=Structure Quat]
struct Quat {
    double x, y, z, w;
    Quat Conjugate() const { return {-x, -y, -z, w}; }
    Quat Inverse()   const { return Conjugate(); }  // si unitaire
    double Norm()    const { return std::sqrt(x*x+y*y+z*z+w*w); }
    Quat Normalized()const { double n=Norm(); return {x/n,y/n,z/n,w/n}; }
    static Quat Identity() { return {0,0,0,1}; }
};
\end{lstlisting}

\begin{center}
\begin{tabular}{ll}
    \toprule
    \textbf{Fonction} & \textbf{Description} \\
    \midrule
    \texttt{FromAxisAngle(n, theta)} & Creation depuis axe-angle \\
    \texttt{Rotate(q, v)} & Rotation d'un vecteur par q \\
    \texttt{ToMat3(q)} & Conversion quaternion vers Mat3d \\
    \texttt{FromMat3(m)} & Conversion matrice vers quaternion \\
    \texttt{Slerp(q1, q2, t)} & Interpolation spherique \\
    \texttt{Lerp(q1, q2, t)} & Interpolation lineaire normalisee \\
    \bottomrule
\end{tabular}
\end{center}

\section{Tests Unitaires (Framework Nkentseu)}

\subsection{Cas de test implementes}

\begin{itemize}
    \item \textbf{TP4\_ConjugateAndNorm} : conjugue et norme d'un quaternion
    \item \textbf{TP4\_Multiplication} : produit de Hamilton
    \item \textbf{TP4\_Rotation} : rotation de $(1,0,0)$ par $\pi/2$ autour de Y
    \item \textbf{TP4\_Slerp} : interpolation entre $q_1$ et $q_2$ a $t=0$ et $t=1$
\end{itemize}

\subsection{Resultats d'execution}

\begin{lstlisting}[caption=Sortie de jenga test -- TP4]
Running test suite: TP4_Tests
Session started: 2026-04-05 16:52:47
Number of tests: 4

[OK]  TP4_ConjugateAndNorm    2/2 assertions  (< 1ms)
[OK]  TP4_Multiplication      1/1 assertions  (< 1ms)
[OK]  TP4_Rotation            2/2 assertions  (< 1ms)
[OK]  TP4_Slerp               2/2 assertions  (< 1ms)

-------------------------------------------------------
RESULTATS DES TESTS
  SUCCES
  Tests :       4 reussis, 4 au total
  Assertions :  7 reussies, 7 au total
  Taux succes : Tests: 100.0%, Assertions: 100.0%
  Temps total : 3ms (< 1ms/test)

Tous les tests sont reussis !
\end{lstlisting}

\section{Code Source}

\begin{center}
\href{https://github.com/neussi/RA_TP4_Quaternions_et_Rotations}{https://github.com/neussi/RA\_TP4\_Quaternions\_et\_Rotations}
\end{center}

\section{Conclusion}

Les quaternions s'averent superieurs aux matrices pour les interpolations de rotations
en animation AR. Le SLERP produit des interpolations a vitesse angulaire constante,
evitant les artefacts visuels des interpolations lineaires.

\end{document}
